\documentclass[12pt,a4paper]{article}

% ── Packages ──────────────────────────────────────────────────────────────────
\usepackage{amsmath, amssymb, bm}
\usepackage{geometry}
\usepackage{booktabs}
\usepackage{array}
\usepackage{xcolor}
\usepackage{hyperref}
\usepackage{parskip}
\usepackage{titlesec}

\geometry{margin=2.5cm}
\hypersetup{colorlinks=true, linkcolor=blue, urlcolor=blue}

% ── Convenience macros ────────────────────────────────────────────────────────
\newcommand{\cph}{\cos\phi}
\newcommand{\sph}{\sin\phi}
\newcommand{\cth}{\cos\theta}
\newcommand{\sth}{\sin\theta}
\newcommand{\cps}{\cos\psi}
\newcommand{\sps}{\sin\psi}
\newcommand{\tph}{\tan\phi}
\newcommand{\tth}{\tan\theta}
\newcommand{\R}{\mathbf{R}}
\newcommand{\q}{\mathbf{q}}
\newcommand{\bu}{\mathbf{u}}
\newcommand{\eone}{\mathbf{e}_1}

\title{\textbf{Robot Kinematics: Motion Models}\\[0.4em]
       \large Differential Drive (2D) \quad$\cdot$\quad
              Frontier Scoring \quad$\cdot$\quad
              Differential Drive (3D)}
\author{}
\date{}

% ══════════════════════════════════════════════════════════════════════════════
\begin{document}
\maketitle
\tableofcontents
\newpage

% ══════════════════════════════════════════════════════════════════════════════
\section{2D Differential Drive Kinematics}
% ══════════════════════════════════════════════════════════════════════════════

\subsection{State Vector and Control Inputs}

The robot's pose is described by three quantities. The control inputs are the
angular velocities of the left and right wheels:

\begin{equation}
  \q(t) = \begin{bmatrix} x(t) \\ y(t) \\ \psi(t) \end{bmatrix},
  \qquad
  \bu(t) = \begin{bmatrix} \omega_L(t) \\ \omega_R(t) \end{bmatrix}
  \tag{1}
\end{equation}

\begin{center}
\begin{tabular}{>{\(}l<{\)} l}
  \toprule
  \text{Symbol} & Meaning \\
  \midrule
  x,\, y        & Position of the robot's centre in the world frame [m] \\
  \psi          & Heading angle measured from the positive $x$-axis [rad] \\
  \omega_L,\,\omega_R & Left and right wheel angular velocities [rad\,s$^{-1}$] \\
  \bottomrule
\end{tabular}
\end{center}

% ──────────────────────────────────────────────────────────────────────────────
\subsection{Wheel Velocities to Linear and Angular Velocity}

Each wheel of radius $r$ contributes a tangential speed. The net linear speed
$v$ and yaw rate $\dot\psi$ are their symmetric and anti-symmetric
combinations, normalised by the wheel-base $L$:

\begin{equation}
  v(t) = \frac{r}{2}\bigl(\omega_R(t) + \omega_L(t)\bigr)
  \tag{2}
\end{equation}

\begin{equation}
  \dot{\psi}(t) = \frac{r}{L}\bigl(\omega_R(t) - \omega_L(t)\bigr)
  \tag{3}
\end{equation}

\begin{center}
\begin{tabular}{>{\(}l<{\)} l}
  \toprule
  \text{Symbol} & Meaning \\
  \midrule
  r & Wheel radius [m] \\
  L & Track width (wheel-to-wheel distance) [m] \\
  v & Forward (linear) speed of the robot [m\,s$^{-1}$] \\
  \dot\psi & Yaw (turning) rate [rad\,s$^{-1}$] \\
  \bottomrule
\end{tabular}
\end{center}

% ──────────────────────────────────────────────────────────────────────────────
\subsection{Continuous-Time Motion Model (ODE)}

Projecting the forward speed onto the world frame gives the full nonlinear
kinematic model:

\begin{equation}
  \dot{\q}(t)
  = \begin{bmatrix} \dot{x}(t) \\ \dot{y}(t) \\ \dot{\psi}(t) \end{bmatrix}
  = \begin{bmatrix} v(t)\cos\psi(t) \\ v(t)\sin\psi(t) \\ \dot{\psi}(t) \end{bmatrix}
  = f\bigl(\q(t),\,\bu(t)\bigr)
  \tag{4}
\end{equation}

% ──────────────────────────────────────────────────────────────────────────────
\subsection{Discrete-Time State Update (Euler Integration)}

For a fixed time-step $\Delta t$, using forward Euler:

\begin{equation}
  \q_{k+1} = \q_k + \Delta t \cdot f(\q_k,\,\bu_k)
  \tag{5}
\end{equation}

Expanding the three state components explicitly:

\begin{equation}
  \begin{aligned}
    x_{k+1}   &= x_k   + v_k \cos\psi_k \cdot \Delta t \\
    y_{k+1}   &= y_k   + v_k \sin\psi_k \cdot \Delta t \\
    \psi_{k+1} &= \psi_k + \dot{\psi}_k  \cdot \Delta t
  \end{aligned}
  \tag{6}
\end{equation}

where $v_k$ and $\dot\psi_k$ are computed from $(\omega_{L,k},\,\omega_{R,k})$
via equations~(2)--(3).

% ──────────────────────────────────────────────────────────────────────────────
\subsection{Exact Arc-Based Integration (Curved Motion)}

When $\dot\psi_k \neq 0$, integrating exactly over the arc avoids Euler
accumulation error. Let $\kappa = \dot\psi_k\,\Delta t$:

\begin{equation}
  \begin{aligned}
    x_{k+1}   &= x_k + \frac{v_k}{\dot{\psi}_k}
                  \Bigl(\sin(\psi_k+\kappa) - \sin\psi_k\Bigr) \\
    y_{k+1}   &= y_k - \frac{v_k}{\dot{\psi}_k}
                  \Bigl(\cos(\psi_k+\kappa) - \cos\psi_k\Bigr) \\
    \psi_{k+1} &= \psi_k + \kappa
  \end{aligned}
  \tag{7}
\end{equation}

When $\dot\psi_k \approx 0$ (straight-line motion), equation~(6) is
numerically preferred; equation~(7) degenerates gracefully via
L'H\^{o}pital's rule to the same result.

% ──────────────────────────────────────────────────────────────────────────────
\subsection{Full Mapping: Control Inputs $\to$ State Update}

Combining equations~(2)--(3) and~(6), the explicit mapping from wheel
velocities to state update is:

\begin{equation}
  \begin{bmatrix} x_{k+1} \\ y_{k+1} \\ \psi_{k+1} \end{bmatrix}
  =
  \begin{bmatrix} x_k \\ y_k \\ \psi_k \end{bmatrix}
  +
  \frac{r\,\Delta t}{2}
  \begin{bmatrix}
    (\omega_{R,k}+\omega_{L,k})\cos\psi_k \\
    (\omega_{R,k}+\omega_{L,k})\sin\psi_k \\
    \dfrac{2}{L}(\omega_{R,k}-\omega_{L,k})
  \end{bmatrix}
  \tag{8}
\end{equation}

% ══════════════════════════════════════════════════════════════════════════════
\newpage
\section{Frontier Scoring Formula}
% ══════════════════════════════════════════════════════════════════════════════

\subsection{The Formula}

Among all candidate frontier regions $R$ on the boundary between known and
unknown space, the robot selects the one with the highest score:

\begin{equation}
  \text{score}(R)
  = \frac{\alpha \cdot \text{size}(R) \;+\; \beta \cdot \text{unknown}(R)}
         {\text{dist}(R) + \varepsilon}
  \tag{9}
\end{equation}

\begin{center}
\begin{tabular}{>{\(}l<{\)} l}
  \toprule
  \text{Symbol} & Meaning \\
  \midrule
  \text{size}(R)    & Number of known free cells in frontier region $R$ \\
  \text{unknown}(R) & Count of unknown cells adjacent to $R$
                      (information gain) \\
  \text{dist}(R)    & Path distance from robot to $R$ [cells or m] \\
  \alpha            & Weight: relative importance of frontier size \\
  \beta             & Weight: relative importance of information gain \\
  \varepsilon       & Small guard constant; prevents division by zero \\
  \bottomrule
\end{tabular}
\end{center}

% ──────────────────────────────────────────────────────────────────────────────
\subsection{Parameter Choices and Rationale}

\paragraph{$\alpha$ — size weight.}
Favours large open frontiers over isolated pockets, preventing the robot from
wasting time on tiny dead-end corridors. Typical range: $\alpha \in [0.3,\,1.0]$.

\paragraph{$\beta$ — unknown weight.}
Directly rewards information gain. Higher $\beta$ drives aggressive
exploration; lower $\beta$ produces cautious, thorough coverage.
Typical range: $\beta \in [1.0,\,3.0]$.

\paragraph{$\varepsilon$ — stability guard.}
Set to a small fraction of the map grid resolution, commonly
$\varepsilon = 10^{-3}$\,m or 1 cell. It must satisfy:
\[
  \varepsilon \;\ll\; \min_R \text{dist}(R)
\]
so that it stabilises the denominator without biasing frontier scores.
A value of $\varepsilon = 10^{-3}$ is appropriate for a 0.05\,m grid.

\paragraph{The $\alpha/\beta$ ratio.}
Scaling both weights by a constant leaves all scores—and therefore the
ranking—unchanged. Only the ratio matters:
\begin{itemize}
  \item $\alpha/\beta < 1$: robot prioritises information density over area.
  \item $\alpha/\beta > 1$: robot prefers large open spaces.
\end{itemize}
A common empirically-tuned starting point is
$\alpha = 0.5,\;\beta = 1.0,\;\varepsilon = 10^{-3}$.

% ──────────────────────────────────────────────────────────────────────────────
\subsection{Failure Modes and Mitigations}

\begin{description}
  \item[Greedy cycling.]
        Two frontiers with equal scores cause oscillation. Mitigation: add a
        hysteresis bonus to the currently-targeted frontier until it is
        resolved.

  \item[Distance dominance late in mapping.]
        When all remaining unknowns are far away, the denominator overwhelms
        the numerator and scores compress. Mitigation: normalise
        $\text{dist}(R)$ by the map diagonal, or use a log penalty
        $\ln(\text{dist}+1)$.

  \item[$\alpha/\beta$ wrong for environment.]
        Open warehouses favour high $\alpha$; tight corridors favour high
        $\beta$. Mitigation: tune offline on a representative map, or adapt
        the ratio based on local occupancy variance.

  \item[$\varepsilon$ too large.]
        If $\varepsilon \gg \text{dist}(R)$ for near frontiers, the distance
        penalty vanishes and nearby frontiers score no better than distant
        ones. Keep $\varepsilon \ll$ minimum expected distance.
\end{description}

% ══════════════════════════════════════════════════════════════════════════════
\newpage
\section{3D Differential Drive Kinematics (Uneven Terrain)}
% ══════════════════════════════════════════════════════════════════════════════

\subsection{Extended State Vector and Control Inputs}

The 6-DOF pose uses ZYX (yaw-pitch-roll) Euler angles. Control inputs remain
the two wheel angular velocities; $z$, roll, and pitch are constrained by the
terrain.

\begin{equation}
  \q(t) = \begin{bmatrix} x(t)\\y(t)\\z(t)\\\phi(t)\\\theta(t)\\\psi(t) \end{bmatrix},
  \qquad
  \bu(t) = \begin{bmatrix}\omega_L(t)\\\omega_R(t)\end{bmatrix}
  \tag{10}
\end{equation}

\begin{center}
\begin{tabular}{>{\(}l<{\)} l}
  \toprule
  \text{Symbol} & Meaning \\
  \midrule
  x,\,y,\,z        & Position in world frame [m] \\
  \phi              & Roll — rotation about world $x$-axis [rad] \\
  \theta            & Pitch — rotation about world $y$-axis [rad] \\
  \psi              & Yaw (heading) — rotation about world $z$-axis [rad] \\
  \omega_L,\,\omega_R & Wheel angular velocities [rad\,s$^{-1}$] \\
  \bottomrule
\end{tabular}
\end{center}

% ──────────────────────────────────────────────────────────────────────────────
\subsection{Step 1 — Wheel Velocities to Body-Frame Speeds}

Identical to the 2D case:

\begin{equation}
  v(t) = \frac{r}{2}\bigl(\omega_R + \omega_L\bigr),
  \qquad
  \dot{\psi}(t) = \frac{r}{L}\bigl(\omega_R - \omega_L\bigr)
  \tag{11}
\end{equation}

% ──────────────────────────────────────────────────────────────────────────────
\subsection{Step 2 — Rotation Matrix $R(\phi,\theta,\psi)$ (ZYX Convention)}

Composing elementary rotations $R_z(\psi)\,R_y(\theta)\,R_x(\phi)$:

\begin{equation}
  R =
  \begin{bmatrix}
    \cps\cth & \cps\sth\sph - \sps\cph & \cps\sth\cph + \sps\sph \\
    \sps\cth & \sps\sth\sph + \cps\cph & \sps\sth\cph - \cps\sph \\
    -\sth    & \cth\sph                & \cth\cph
  \end{bmatrix}
  \tag{12}
\end{equation}

where $c_\alpha \equiv \cos\alpha$, $s_\alpha \equiv \sin\alpha$.
The first column of $R$ is the forward (body $x$) direction in world
coordinates.

% ──────────────────────────────────────────────────────────────────────────────
\subsection{Step 3 — World-Frame Translational Velocity}

The robot moves at speed $v$ along its body $x$-axis. Multiplying by the
first column of $R$:

\begin{equation}
  \begin{bmatrix}\dot x \\ \dot y \\ \dot z\end{bmatrix}
  = v(t)\cdot R(\phi,\theta,\psi)\,\eone
  = v(t)\begin{bmatrix}\cps\cth \\ \sps\cth \\ -\sth\end{bmatrix}
  \tag{13}
\end{equation}

\textit{Note:} On flat ground ($\theta=0$), this reduces exactly to the 2D
equations. On an uphill slope ($\theta>0$), $\dot z = -v\sin\theta > 0$.

% ──────────────────────────────────────────────────────────────────────────────
\subsection{Step 4 — Euler Kinematic Equations (Angular Velocity)}

The Euler angle rates and body angular velocity
$\bm\omega_b = (p,\,q,\,r)^\top$ are related by:

\begin{equation}
  \begin{bmatrix}\dot\phi\\\dot\theta\\\dot\psi\end{bmatrix}
  =
  \underbrace{
  \begin{bmatrix}
    1 & \sph\tth & \cph\tth \\
    0 & \cph     & -\sph    \\
    0 & \sph/\cth & \cph/\cth
  \end{bmatrix}}_{T(\phi,\theta)^{-1}}
  \begin{bmatrix}p\\q\\r\end{bmatrix}
  \tag{14}
\end{equation}

\textit{Singularity:} $T^{-1}$ is undefined at $\theta = \pm90°$.
For steep terrain, replace Euler angles with a unit quaternion
$\mathbf{q} = (q_0,q_1,q_2,q_3)^\top$ and use:
\[
  \dot{\mathbf{q}} = \tfrac{1}{2}\,\Xi(\mathbf{q})\,\bm\omega_b
\]

For a ground robot, the body angular velocity components are:

\begin{equation}
  p = \dot\phi_{\text{surface}}, \quad
  q = \dot\theta_{\text{surface}}, \quad
  r = \dot\psi
  \tag{15}
\end{equation}

where $\dot\phi_\text{surface}$ and $\dot\theta_\text{surface}$ are the
time-derivatives of the terrain surface normal projected onto the robot body,
and $\dot\psi$ comes from equation~(11).

% ──────────────────────────────────────────────────────────────────────────────
\subsection{Step 5 — Surface Constraint (Terrain Model)}

On a surface $z = h(x,y)$ with gradient
$\nabla h = (\partial h/\partial x,\;\partial h/\partial y)$:

\begin{align}
  \theta(t) &= -\arctan\!\left(
      \frac{\partial h}{\partial x}\cos\psi
    + \frac{\partial h}{\partial y}\sin\psi
    \right)
  \tag{16a} \\[6pt]
  \phi(t)   &= \arctan\!\left(
    \frac{
      -\dfrac{\partial h}{\partial x}\sin\psi
      + \dfrac{\partial h}{\partial y}\cos\psi
    }{
      \sqrt{1 + \Bigl(
        \dfrac{\partial h}{\partial x}\cos\psi
        + \dfrac{\partial h}{\partial y}\sin\psi
      \Bigr)^2}
    }
    \right)
  \tag{16b} \\[6pt]
  z(t) &= h\!\bigl(x(t),\,y(t)\bigr)
  \tag{16c}
\end{align}

Equations~(16a--c) mean $z$, $\phi$, $\theta$ are \emph{algebraically
constrained} by the terrain at each step; only $x$, $y$, $\psi$ are
independently integrated.

% ──────────────────────────────────────────────────────────────────────────────
\subsection{Complete Discrete-Time State Update (Euler Integration)}

Combining all steps over timestep $\Delta t$:

\begin{equation}
  \boxed{
  \begin{aligned}
    x_{k+1}      &= x_k      + v_k\,\cps_k\cth_k\,\Delta t \\
    y_{k+1}      &= y_k      + v_k\,\sps_k\cth_k\,\Delta t \\
    z_{k+1}      &= h(x_{k+1},\,y_{k+1}) \\
    \phi_{k+1}   &= \phi_{\text{surface}}(x_{k+1},\,y_{k+1},\,\psi_{k+1}) \\
    \theta_{k+1} &= \theta_{\text{surface}}(x_{k+1},\,y_{k+1},\,\psi_{k+1}) \\
    \psi_{k+1}   &= \psi_k   + \dot\psi_k\,\Delta t
  \end{aligned}
  }
  \tag{17}
\end{equation}

Only $x$, $y$, $\psi$ are numerically integrated; the remaining three states
are read off the terrain at each step via equations~(16a--c).

% ──────────────────────────────────────────────────────────────────────────────
\subsection{Unconstrained 3D Extension (Free Body — Drone / AUV)}

If the robot is not surface-constrained (aerial, underwater), all six states
must be integrated with independent actuation on all axes:

\begin{equation}
  \begin{bmatrix}\dot x\\\dot y\\\dot z\end{bmatrix}
  = R(\phi,\theta,\psi)
    \begin{bmatrix}v_x^b\\v_y^b\\v_z^b\end{bmatrix},
  \qquad
  \begin{bmatrix}\dot\phi\\\dot\theta\\\dot\psi\end{bmatrix}
  = T(\phi,\theta)^{-1}
    \begin{bmatrix}p\\q\\r\end{bmatrix}
  \tag{18}
\end{equation}

All six components $(p,q,r,v_x^b,v_y^b,v_z^b)$ are independent inputs with no
wheel coupling constraint.

% ══════════════════════════════════════════════════════════════════════════════
\end{document}
